% ============================================================================
\section{Introduction}
\label{sec:intro}
% ============================================================================
Identifying tau particles is important in studies of particles such as the
Higgs and Z bosons. Many taus decay into a narrow spray of particles, which
leaves a compact energy deposit in the ATLAS calorimeter. Ordinary jets of
particles can leave a similar deposit, especially at low energy. A simple
energy threshold therefore either accepts too much background or rejects useful
tau candidates.

The decision must be made while collision data are arriving. The Large Hadron
Collider (LHC) produces proton collisions at \SI{40}{MHz}, far faster than
ATLAS can store them. Its first-stage, or Level-1, trigger therefore keeps only
a small fraction for later reconstruction~\cite{atlas_trigger_run3,atlas_phase1_tdr}.
We study compact neural networks for the calorimeter Trigger Objects (TOBs)
supplied by the electron Feature Extractor (eFEX). The comparison is made at a
fixed event fake rate of $0.5\%$: a useful model must retain more signal without
accepting more background events.

We first study models that score one TOB at a time. We test which detector
inputs carry useful information, whether compact convolutional networks improve
on fully connected networks, and how the classifier and loss can protect the required efficiency
across the energy range. We then ask a different question: can the network make
a better event decision if it sees two TOBs together? The pair model applies
the same small network to both TOBs, scores every pair in the event and uses the
highest pair score. During training, every event receives equal weight even
when it contains many pairs.

For the single-TOB study, the strongest global model accepts $69.61\%$ of
truth-$\tau$ objects, compared with $57.37\%$ for the eFEX baseline. The
pair study uses a different measure: the fraction of signal events that fire.
On its final test, this efficiency rises from $34.30\%$ for the second-highest-
$\ptTOB$ baseline to $46.70\%$, a gain of 12.40 percentage points. The three
models reach event fake rates of $0.5183$--$0.5223\%$, slightly above the
$0.5\%$ limit. The efficiency improvement is therefore clear, but the final
success criterion is not met. The two studies use different denominators and
their percentages should not be compared directly.

\noindent\textbf{Code and demo:} \url{\repourl}
